% GENERATED FILE. Edit canonical lessons, then run npm run generate:books.
% canonical-source-hash: fnv1a64:49f635e1f6fe24b1
% canonical-lessons: KA-C494-mududu, KA-C494-tadavarisu, KA-C494-beragu, KA-C494-kangedu, KA-C494-digbhrame

\chapter{Shock and Dismay}
\label{ch:ka-shock-and-dismay}
\hlchaptermodality{494}

\begin{chapteropening}
\textbf{By the end of this chapter:} \emph{I can say to shrink back, to fumble, amazement, to lose heart and bewilderment.}

Complete the last lesson of chapter 494: I can say to shrink back, to fumble, amazement, to lose heart and bewilderment.
\end{chapteropening}

\section[\texorpdfstring{\kn{ಮುದುಡು} (muduḍu)}{ಮುದುಡು (muduḍu)}]{\kn{ಮುದುಡು} (muduḍu) — to shrink back; to curl up}

\label{lesson:KA-C494-mududu}



\begin{quote}
Before the new one: say the Kannada for frustration, then the Kannada for agitation.
\end{quote}

\subsection*{You'll want to know: \kn{ಮುದುಡು}}
\textbf{\kn{ಮುದುಡು}} — \emph{muduḍu} — \textquotedblleft{}to shrink back; to curl up\textquotedblright{}.

\begin{quote}
\kn{ಗುಡುಗಿನ} \kn{ಸದ್ದಿಗೆ} \kn{ನಾಯಿ} \kn{ಮೂಲೆಯಲ್ಲಿ} \kn{ಮುದುಡಿ} \kn{ಮಲಗಿತು}.
\end{quote}

\begin{quote}
\emph{guḍugina saddige nāyi mūleyalli muduḍi malagitu.}
\end{quote}

\begin{quote}
At the sound of thunder the dog curled up in a corner.
\end{quote}

Past (I …): \kn{ಮುದುಡಿದೆ}, \emph{muduḍide}.

Of fear, cold or a wilting leaf; \kn{ಮುದುಡಿಕೊ}, to huddle up.

\begin{cousinweb}
\textbf{\kn{ಮುದುಡು}} (\emph{muduḍu}) is a native Kannada verb, \textquotedblleft{}to shrink, to curl up\textquotedblright{}; its origin is not settled.
\end{cousinweb}

\begin{grammarlens}[title={What you've built}]
One more word for describing what you lived and felt.
\end{grammarlens}

\subsection*{Your turn}
\begin{itemize}
\raggedright
  \item \emph{Say it:} \emph{muduḍu}
  \item \emph{Say it:} \emph{muduḍu}, once more
  \item \emph{Say it:} \emph{guḍugina saddige nāyi mūleyalli muduḍi malagitu.} — At the sound of thunder the dog curled up in a corner.
  \item \emph{Recall:} say the Kannada for frustration, then the Kannada for agitation, then say \emph{muduḍu} again
\end{itemize}

\begin{tcolorbox}[breakable,colback=teal!4,colframe=teal!35!black,title={Before you move on}]
What does \kn{ಮುದುಡು} mean? (\textquotedblleft{}To shrink back; to curl up\textquotedblright{}.) Say it once more.
\end{tcolorbox}
\section[\texorpdfstring{\kn{ತಡವರಿಸು} (taḍavarisu)}{ತಡವರಿಸು (taḍavarisu)}]{\kn{ತಡವರಿಸು} (taḍavarisu) — to fumble; to falter}

\label{lesson:KA-C494-tadavarisu}



\begin{quote}
Before the new one: say the Kannada for agitation, then the Kannada for to shrink back.
\end{quote}

\subsection*{You'll want to know: \kn{ತಡವರಿಸು}}
\textbf{\kn{ತಡವರಿಸು}} — \emph{taḍavarisu} — \textquotedblleft{}to fumble; to falter\textquotedblright{}.

\begin{quote}
\kn{ಪ್ರಶ್ನೆ} \kn{ಕೇಳಿದಾಗ} \kn{ಅವನು} \kn{ಉತ್ತರಕ್ಕಾಗಿ} \kn{ತಡವರಿಸಿದನು}.
\end{quote}

\begin{quote}
\emph{praśne kēḷidāga avanu uttarakkāgi taḍavarisidanu.}
\end{quote}

\begin{quote}
When he was asked, he fumbled for an answer.
\end{quote}

Past (I …): \kn{ತಡವರಿಸಿದೆ}, \emph{taḍavariside}.

For words: \kn{ಮಾತಿಗೆ} \kn{ತಡವರಿಸು}, to grope for words. In the dark it is feeling your way.

\begin{cousinweb}
\textbf{\kn{ತಡವರಿಸು}} (\emph{taḍavarisu}) is native Dravidian: DEDR 3025 gives Kannada \emph{taḍavu}, \textquotedblleft{}to touch gently, grope\textquotedblright{}, and \emph{taḍa-varisu}, \textquotedblleft{}to feel one's way with the hands or feet\textquotedblright{}.
\end{cousinweb}

\begin{grammarlens}[title={What you've built}]
One more word for describing what you lived and felt.
\end{grammarlens}

\subsection*{Your turn}
\begin{itemize}
\raggedright
  \item \emph{Say it:} \emph{taḍavarisu}
  \item \emph{Say it:} \emph{taḍavarisu}, once more
  \item \emph{Say it:} \emph{praśne kēḷidāga avanu uttarakkāgi taḍavarisidanu.} — When he was asked, he fumbled for an answer.
  \item \emph{Recall:} say the Kannada for agitation, then the Kannada for to shrink back, then say \emph{taḍavarisu} again
\end{itemize}

\begin{tcolorbox}[breakable,colback=teal!4,colframe=teal!35!black,title={Before you move on}]
What does \kn{ತಡವರಿಸು} mean? (\textquotedblleft{}To fumble; to falter\textquotedblright{}.) Say it once more.
\end{tcolorbox}
\section[\texorpdfstring{\kn{ಬೆರಗು} (beragu)}{ಬೆರಗು (beragu)}]{\kn{ಬೆರಗು} (beragu) — amazement; astonishment}

\label{lesson:KA-C494-beragu}



\begin{quote}
Before the new one: say the Kannada for to shrink back, then the Kannada for to fumble.
\end{quote}

\subsection*{You'll want to know: \kn{ಬೆರಗು}}
\textbf{\kn{ಬೆರಗು}} — \emph{beragu} — \textquotedblleft{}amazement; astonishment\textquotedblright{}.

\begin{quote}
\kn{ಜಲಪಾತವನ್ನು} \kn{ನೋಡಿ} \kn{ಮಕ್ಕಳು} \kn{ಬೆರಗಾದರು}.
\end{quote}

\begin{quote}
\emph{jalapātavannu nōḍi makkaḷu beragādaru.}
\end{quote}

\begin{quote}
Seeing the waterfall, the children were amazed.
\end{quote}

\kn{ಬೆರಗಾಗು}, to be amazed: the noun joined to \kn{ಆಗು}.

\begin{cousinweb}
\textbf{\kn{ಬೆರಗು}} (\emph{beragu}) is native Dravidian: DEDR 5443 lists Kannada \emph{beṟagu}, \textquotedblleft{}amazement, astonishment, alarm\textquotedblright{}, beside Tamil words for being bewildered.
\end{cousinweb}

\begin{grammarlens}[title={What you've built}]
One more word for describing what you lived and felt.
\end{grammarlens}

\subsection*{Your turn}
\begin{itemize}
\raggedright
  \item \emph{Say it:} \emph{beragu}
  \item \emph{Say it:} \emph{beragu}, once more
  \item \emph{Say it:} \emph{jalapātavannu nōḍi makkaḷu beragādaru.} — Seeing the waterfall, the children were amazed.
  \item \emph{Recall:} say the Kannada for to shrink back, then the Kannada for to fumble, then say \emph{beragu} again
\end{itemize}

\begin{tcolorbox}[breakable,colback=teal!4,colframe=teal!35!black,title={Before you move on}]
What does \kn{ಬೆರಗು} mean? (\textquotedblleft{}Amazement; astonishment\textquotedblright{}.) Say it once more.
\end{tcolorbox}
\section[\texorpdfstring{\kn{ಕಂಗೆಡು} (kaṅgeḍu)}{ಕಂಗೆಡು (kaṅgeḍu)}]{\kn{ಕಂಗೆಡು} (kaṅgeḍu) — to lose heart; to be dismayed}

\label{lesson:KA-C494-kangedu}



\begin{quote}
Before the new one: say the Kannada for to fumble, then the Kannada for amazement.
\end{quote}

\subsection*{You'll want to know: \kn{ಕಂಗೆಡು}}
\textbf{\kn{ಕಂಗೆಡು}} — \emph{kaṅgeḍu} — \textquotedblleft{}to lose heart; to be dismayed\textquotedblright{}.

\begin{quote}
\kn{ಪರೀಕ್ಷೆಯಲ್ಲಿ} \kn{ಫೇಲಾದರೂ} \kn{ಅವನು} \kn{ಕಂಗೆಡಲಿಲ್ಲ}.
\end{quote}

\begin{quote}
\emph{parīkṣeyalli phēlādarū avanu kaṅgeḍalilla.}
\end{quote}

\begin{quote}
Though he failed the exam, he did not lose heart.
\end{quote}

Past (I …): \kn{ಕಂಗೆಟ್ಟೆ}, \emph{kaṅgeṭṭe}.

Often in the negative, as praise: \kn{ಕಂಗೆಡಲಿಲ್ಲ}, did not lose heart.

\begin{cousinweb}
\textbf{\kn{ಕಂಗೆಡು}} (\emph{kaṅgeḍu}) is a native Kannada verb, \textquotedblleft{}to be dismayed, to lose heart\textquotedblright{}; its origin is not settled.
\end{cousinweb}

\begin{grammarlens}[title={What you've built}]
One more word for describing what you lived and felt.
\end{grammarlens}

\subsection*{Your turn}
\begin{itemize}
\raggedright
  \item \emph{Say it:} \emph{kaṅgeḍu}
  \item \emph{Say it:} \emph{kaṅgeḍu}, once more
  \item \emph{Say it:} \emph{parīkṣeyalli phēlādarū avanu kaṅgeḍalilla.} — Though he failed the exam, he did not lose heart.
  \item \emph{Recall:} say the Kannada for to fumble, then the Kannada for amazement, then say \emph{kaṅgeḍu} again
\end{itemize}

\begin{tcolorbox}[breakable,colback=teal!4,colframe=teal!35!black,title={Before you move on}]
What does \kn{ಕಂಗೆಡು} mean? (\textquotedblleft{}To lose heart; to be dismayed\textquotedblright{}.) Say it once more.
\end{tcolorbox}
\section[\texorpdfstring{\kn{ದಿಗ್ಭ್ರಮೆ} (digbhrame)}{ದಿಗ್ಭ್ರಮೆ (digbhrame)}]{\kn{ದಿಗ್ಭ್ರಮೆ} (digbhrame) — bewilderment; stunned shock}

\label{lesson:KA-C494-digbhrame}



\begin{quote}
Before the new one: say the Kannada for amazement, then the Kannada for to lose heart.
\end{quote}

\subsection*{You'll want to know: \kn{ದಿಗ್ಭ್ರಮೆ}}
\textbf{\kn{ದಿಗ್ಭ್ರಮೆ}} — \emph{digbhrame} — \textquotedblleft{}bewilderment; stunned shock\textquotedblright{}.

\begin{quote}
\kn{ಮನೆಗೆ} \kn{ಬಂದಾಗ} \kn{ಬಾಗಿಲು} \kn{ತೆರೆದಿತ್ತು}, \kn{ನಾನು} \kn{ದಿಗ್ಭ್ರಮೆಗೊಂಡೆ}.
\end{quote}

\begin{quote}
\emph{manege bandāga bāgilu teredittu, nānu digbhramegoṇḍe.}
\end{quote}

\begin{quote}
When I came home the door was open; I was stunned.
\end{quote}

\kn{ದಿಗ್ಭ್ರಮೆಗೊಳ್ಳು}, to be stunned. Literally a whirling of the directions: not knowing which way to turn.

\begin{cousinweb}
\textbf{\kn{ದಿಗ್ಭ್ರಮೆ}} (\emph{digbhrame}) is borrowed from Sanskrit \emph{digbhrama}, \textquotedblleft{}perplexity about the points of the compass\textquotedblright{} (Monier-Williams): \emph{dik}, \textquotedblleft{}direction\textquotedblright{}, and \emph{bhrama}, \textquotedblleft{}giddiness\textquotedblright{} (CDIAL 9646).
\end{cousinweb}

\begin{grammarlens}[title={What you've built}]
One more word for describing what you lived and felt.
\end{grammarlens}

\subsection*{Your turn}
\begin{itemize}
\raggedright
  \item \emph{Say it:} \emph{digbhrame}
  \item \emph{Say it:} \emph{digbhrame}, once more
  \item \emph{Say it:} \emph{manege bandāga bāgilu teredittu, nānu digbhramegoṇḍe.} — When I came home the door was open; I was stunned.
  \item \emph{Recall:} say the Kannada for amazement, then the Kannada for to lose heart, then say \emph{digbhrame} again
\end{itemize}

\begin{tcolorbox}[breakable,colback=teal!4,colframe=teal!35!black,title={Before you move on}]
What does \kn{ದಿಗ್ಭ್ರಮೆ} mean? (\textquotedblleft{}Bewilderment; stunned shock\textquotedblright{}.) Say it once more.
\end{tcolorbox}
